% Lösungen Einheit 1 von 4 – Punktprobe und Seitenlage (zusammenbau v0.1)
\einheitenkopf{Einheit 1 von 4 · Punktprobe und Seitenlage}

\erg{9}{a) Punkt gegen Ebene – die Prüfregel ist die Gleichung von $E$.}
\erg{10}{a) auf der Seite mit größerem Wert – $6 > 3$. \quad b) $2 \cdot 3 - 1 + 4 = 9$, und $9 \neq 7$ – $P$ liegt nicht in $E$. \quad c) $2 - 4 \cdot 2 = -6$ – erfüllt, $S$ liegt in $E$; die x-Koordinate kommt in der Gleichung nicht vor und darf beliebig sein. \quad d) aus x: $t = 2$, aus z: $s = 3$; y: $0 + 2 + 3 = 5$ – die dritte Gleichung stimmt, $P$ liegt in $E$. \quad e) $5 \cdot 2 - 2 \cdot 5 = 0$ – $K$ liegt in $W$; $(5 | -2 | 0) \cdot (0 | 0 | 1) = 0$ – die Normalenvektoren stehen senkrecht aufeinander, also steht $W$ senkrecht auf der Fahrbahn. \quad f) „Untersuche, ob die Ebene $W$ die Kante $AE$ schneidet, und gib gegebenenfalls den Punkt an.“ Aus (2): $r = 2$, im Bereich – $W$ schneidet die Kante in $(5 | 0 | 2)$. \quad g) Eingesetzt: $1 + 2 \cdot 0,5 + 1 = 3$; wegen $2 < 3 < 6$ liegt $T$ bei $L_1$ auf der vom Ursprung abgewandten Seite und bei $L_2$ auf der Ursprungsseite; alle Koordinaten sind positiv – $T$ liegt im Inneren. \quad h) z. B. $P(1 | 2 | 4)$: $2 + 2 + 4 = 8$ und alle Koordinaten positiv – $P$ liegt im Innern; $Q(2 | 1 | 4)$: $4 + 1 + 4 = 9 \neq 8$ – $Q$ liegt nicht einmal in $E$.}
\erg{11}{a) Lea hat die dritte Gleichung nicht geprüft: y: $0 + 2 \cdot 2 + 2 = 6$, und $6 \neq 5$ – $P$ liegt nicht in $E$. \quad b) Die Ebene ist genau die Menge der Punkte, deren Koordinaten die Gleichung erfüllen – erfüllt heißt drin, nicht erfüllt heißt draußen.}
